\subsection{Implementation and deployment}\label{sec:implementation}

Both programs are written in C99 and built by one GNU make file; they need a C99
compiler (GCC~4.7 or later, or Clang), libjpeg-turbo~1.4 or later and, for \hellion{},
FFmpeg at run time (not on a vehicle that sends its camera's own H.264), and they run
on Linux~3.2 or later. Release binaries are built
statically for x86-64, ARMv7 with NEON and AArch64 with GCC~15.3 (Bootlin toolchains
2026.08-1 \cite{bootlin}) and libjpeg-turbo~3.2.0, so that they need no libraries
on the target. Child processes are started with their program resolved before
\texttt{fork}, with inherited file descriptors closed and signals reset, in their own
process group, and are stopped by closing their input before signalling them. An
installation script installs both programs as system services on each end of the
tether, given the end's role and the placement and, for a camera that encodes, its device,
in which case FFmpeg is neither installed nor required on the vehicle; the two services run as one service
user, so that \hellion{} can write \eeye{}'s configuration and both can use the frame
bus, and neither needs elevated privileges.
